%!TEX root=main.tex

\section{Introduction}\label{sect:intro}
\EZComm{DA FARE, intanto ho messo questo pezzettino con lo schema dei tre step}
We illustrate our approach on a tiny imperative Java-like language, in the spirit of the seminal calculus Featherweight Java (FJ) \cite{IgarashiPW99}. 
Besides the FJ constructs, the language includes some imperative ones, notably field update and deallocation, and evaluation takes place in a memory, a map from references (unique object identifiers) into object states. 
\EZComm{ci sono stati diversi calcoli imperativi in letteratura forse citarli}

To modularly present the key features of our approach, we consider three  versions sof increasing complexity. 

In \cref{sect:flat}, field types can only be primitive, hence, in the object state associated to a reference, field values cannot be references in turn.
In other words, the memory is \emph{flat}, and references can only be accessed directly through variables. 


In \cref{sect:hierarchical}, we remove the restriction on field types, hence references can be accessed through \emph{paths}, e.g., $\FieldAccess{\FieldAccess{\x}{\f_1}\ldots}{\f_n}$. 
Hence the memory is \emph{hierarchical}. 

\bez
Finally, in \cref{sect:promotion} we 
enrich the type system by keeping trace of the aliasing introduced among paths. In this way, typechecking a sequential composition can incorporate  a \emph{promotion} \EZComm{recovery?} mechanism, allowing to assign more restrictive assumptions to variables in the second subterm. 
Notably, in the context used to typecheck the first subterm we distinguish the variables which are \emph{borrowed}, that is, cannot be accessed through an alias in the second subterm.
\eez

